\section{Results}
\label{sec:results}

Every number in this section is produced by an analysis and inserted by a macro from the frozen result
files. A hypothesis is named by its number in the register (Appendix~\ref{app:register}), which gives its
rule, its status and whether it was declared before or after earlier results were seen; the status quoted
here is the one computed from the stored results. Fractions are over subjects with Wilson intervals;
paired differences are over subjects with bootstrap intervals; ``interior'' is defined in
Section~\ref{sec:estimation}.

\subsection{Structural identifiability depends on the gut's initial state, and the rank is full}
\label{sec:structural}

Symbolic analysis with StructuralIdentifiability.jl (Appendix~\ref{app:structural}) asks what is determined
by an ideal observation of glucose. The answer depends on an assumption that the data cannot settle:
what is known about the gut at the start of a meal (Table~\ref{tab:structural}). When the intestinal
content is a generic, known value, the tool reports $\ke$ and $\ka$ globally identifiable, which is the
formulation fixed in the plan. When the gut state is unconstrained or unknown, they are only locally
identifiable. In canonical coordinates the two symmetric functions $\ke+\ka$ and $\ke\ka$ are globally
identifiable, which is the unordered pair. The simulator starts each meal with an empty gut, the case of
Proposition~\ref{prop:exchange}, where the exchange is an exact symmetry; numerically, exchanging the rates
changes simulated glucose by at most $\resGutSweepMaxZero$\,mg\,dL$^{-1}$.

\begin{table}[t]
\centering\small
\begin{tabularx}{\linewidth}{Xccc}
\toprule
formulation & initial gut state & $\SI$ & $\ke,\ka$ \\
\midrule
meal as input, all initial conditions generic and known & known, generic & global & global \\
\emph{planned}: meal as initial stomach content, no input & known, generic & global & global \\
meal as input & unconstrained or unknown & global & local \\
canonical coordinates $(\ke+\ka,\ \ke\ka)$ & unconstrained or unknown & global & the two combinations are global \\
\bottomrule
\end{tabularx}
\caption{Structural identifiability of the model with glucose as output, by formulation. A generic known
intestinal content makes the rates globally identifiable; an unconstrained one leaves the unordered pair.}
\label{tab:structural}
\end{table}

The plan's rule was to delete the exchange framing if the rates were globally identifiable. It fired for
the planned formulation, so H1 is reported as \emph{formulation-dependent}, not as met
(\resVerdictHOne). We keep the exchange framing because the empty-gut model that the simulator runs, and
that free-living data approximate to the extent that meals are isolated, is exactly symmetric.

The stacked sensitivity matrix has generic rank three at relative tolerance $10^{-6}$ in
\resRankFullIauc{} of \resRankDrawsIauc{} random parameter draws for iAUC, and in every draw for iAUC with
centroid and for the trace. The three parameters are therefore locally identifiable under every
observable of the ladder; what differs is conditioning, with the ratio of the smallest to the largest
singular value \resRankRatioIauc{} for iAUC, \resRankRatioCentroid{} with the centroid and
\resRankRatioTrace{} for the trace. We accordingly describe iAUC as locally identifiable but
ill-conditioned, and make no stronger structural claim than the exchange symmetry.

\subsection{$\ke$ and $\ka$ cannot be separated from the area, in any cohort}
\label{sec:timing}

\paragraph{The estimates sit on their bounds.}
At the primary box the maximum-likelihood estimates of $\ke$ and $\ka$ are on a bound in \resPinnedFracKeOne\%
and \resPinnedFracKaOne\% of the \resSubjectsOne{} CGMacros subjects, on the upper bound in
\resPinnedUpperKeOne{} and \resPinnedUpperKaOne{}; those of $\SI$ in \resPinnedFracSIOne\%. Widening the box
to twice its width lowers the fractions for $\ke$ and $\ka$ to \resPinnedFracKeTwo\% and
\resPinnedFracKaTwo\%, and halving it raises them to \resPinnedFracKeHalf\% and \resPinnedFracKaHalf\%;
they do not reach zero (Appendix, Figure~\ref{fig:gutsweep}b). The estimates are not an artefact of the
optimizer: the negative log-likelihood of the Adam estimate exceeds that of the best of six L-BFGS runs by
at most \resOptGapMaxHalf, \resOptGapMaxOne{} and \resOptGapMaxTwo{} at the three widths, always below
$\delta$.

\paragraph{Fisher information.}
In log coordinates the median eigenvalues of the information matrix at the primary box are \resEigMedianOne,
\resEigMedianTwo{} and \resEigMedianThree{} (Figure~\ref{fig:fisher}a). The weakest eigenvector of the
timing block lies along the exchange direction, with a median absolute cosine of \resCosineIaucOne{} under
iAUC, \resCosineCentroidOne{} with the centroid and \resCosineTraceOne{} under the trace (H4,
\resVerdictHFour); Figure~\ref{fig:symmetry}b shows that this holds for every subject in all three cohorts.

\begin{figure}[t]
\centering
\safefigure[width=\linewidth]{f1_symmetry}
\caption{\textbf{(a)} Rate of gut appearance for a meal at two rate pairs related by exchange, in an empty
gut (the curves coincide) and with $5\%$ of the meal left in the intestine (dotted; the symmetry is
broken). This panel is analytic, not data. \textbf{(b)} Cumulative distribution over subjects of
$1-|\cos|$ between the weakest Fisher direction of the timing block and the exchange direction
$(1/\ka,-1/\ke)$, iAUC objective, primary box. The dashed line marks the threshold of H4 and H15
($|\cos|=0.8$); every subject of every cohort lies far on the near side of it.}
\label{fig:symmetry}
\end{figure}

\paragraph{Profile likelihood.}
No profile interval for $\ke$ or $\ka$ is bounded under iAUC: \resProfileIaucKeOne\% and
\resProfileIaucKaOne\% of subjects at the primary box (upper Wilson limits \resProfileIaucKeOneHigh\% and
\resProfileIaucKaOneHigh\%), and the same at half and twice the box. Extending the grid to four times the
upper bound leaves \resExtFlatKeOne{} and \resExtFlatKaOne{} of \resSubjectsOne{} profiles flat and the
rest one-sided (Figure~\ref{fig:profiles}). The clause of H5 for the timing parameters under iAUC, at most
a quarter of subjects bounded, therefore holds. This is not a statement that the data are uninformative: the
information is in a combination of the two rates (Section~\ref{sec:more}), and in the area deficit below.

\begin{figure}[t]
\centering
\safefigure[width=\linewidth]{f3_profile_curves}
\caption{Profile negative log-likelihood above its minimum for $\SI$, $\ke$ and $\ka$ (columns) for two
subjects (rows), iAUC objective, primary box. Dotted: $\delta$. Grey dashed: the extension to four times
the upper bound. \textbf{(a)} A subject whose $\SI$ profile is steep on one side and bounded by the box
on the other. \textbf{(b)} A subject whose timing profiles never reach $\delta$. Neither $\ke$ nor $\ka$ is
bounded in the first subject either.}
\label{fig:profiles}
\end{figure}

\paragraph{Pinning is partly an area deficit.}
Subjects whose $\ke$ or $\ka$ sits on the upper bound are not simply those with a flat likelihood. They
have much larger responses (mean observed iAUC \resBiasObsPinnedKe{} against \resBiasObsOtherKe{} for the
others) and are under-predicted by \resBiasSignedKe{} mg\,dL$^{-1}$\,min on average, against
\resBiasSignedOtherKe{} for the others: the model cannot deliver the observed area at any admissible
timing, and the fit pushes the rates to their limit. Part of the pinning is therefore misfit, which is a
finding about the model and the carbohydrate logs, not about identifiability, and which the replica of
Section~\ref{sec:si} lets us separate.

\paragraph{Three cohorts.}
The same picture holds in the Shanghai cohort (\resRepShanghaiOneN{} subjects, $15$-minute sensor, dietary
records) and in the cohort with standardized meals (\resRepHallOneN{} subjects)
(Table~\ref{tab:cohortresults}, Figure~\ref{fig:cohorts}). No interval for $\ke$ or $\ka$ is bounded in
either cohort at the primary or the double box, and the weakest Fisher direction has $|\cos|\ge0.8$
with the exchange direction in \resRepShanghaiOneCosine\% and \resRepHallOneCosine\% of subjects (H15,
\resVerdictHFifteen). We note a limit of the evidence: at the primary box most estimates in these cohorts
are on a bound (Table~\ref{tab:cohortresults}), so only \resRepShanghaiOneInteriorKe{} and
\resRepShanghaiOneInteriorKa{} Shanghai subjects and \resRepHallOneInteriorKe{} and
\resRepHallOneInteriorKa{} subjects of the other cohort are interior for $\ke$ and $\ka$; the
replication rests mainly on the Fisher alignment and the flat profiles, not on the interior fractions.

\begin{table}[t]
\centering\small
\begin{tabular}{llrcccc}
\toprule
cohort & box & subjects & estimate on a bound (\%) & bounded interval (\%) & interior for $\ke$ / $\ka$ & $|\cos|\ge0.8$ (\%) \\
 & & & $\SI$ / $\ke$ / $\ka$ & $\SI$ / $\ke$ / $\ka$ & & \\
\midrule
CGMacros & $1\times$ & 45 & 51 / 89 / 87 & 7 / 0 / 0 & 5 / 6 & 100 \\
CGMacros & $2\times$ & 45 & 11 / 58 / 76 & 53 / 0 / 0 & 19 / 11 & 100 \\
Shanghai & $1\times$ & 87 & 45 / 95 / 97 & 7 / 0 / 0 & 4 / 2 & 100 \\
Shanghai & $2\times$ & 87 & 10 / 68 / 72 & 45 / 0 / 0 & 24 / 24 & 100 \\
Standardized meals & $1\times$ & 15 & 73 / 80 / 67 & 0 / 0 / 0 & 3 / 5 & 100 \\
Standardized meals & $2\times$ & 15 & 7 / 60 / 40 & 0 / 0 / 0 & 6 / 9 & 100 \\
\bottomrule
\end{tabular}
\caption{Replication of the identifiability result on three cohorts, iAUC objective. The last column is the share of subjects whose weakest timing Fisher direction has an absolute cosine of at least $0.8$ with the exchange direction. Interior counts are subjects whose estimate of the parameter is more than one percent of the box width from both bounds.}
\label{tab:cohortresults}
\end{table}


\begin{figure}[t]
\centering
\safefigure[width=\linewidth]{f6_cohorts}
\caption{Replication across cohorts, iAUC objective, primary box. \textbf{(a)} Subjects whose estimate of a
parameter lies on a bound. \textbf{(b)} Subjects with a bounded $95\%$ profile interval, with Wilson
intervals. No interval for $\ke$ or $\ka$ is bounded in any cohort.}
\label{fig:cohorts}
\end{figure}

\subsection{$\SI$ is recoverable for some subjects and not others, and not mainly because of noise}
\label{sec:si}

\paragraph{Real data.}
The plan required the $\SI$ interval to be bounded in at least $80\%$ of subjects interior for $\SI$. It is
bounded in \resProfileIaucSIInteriorBoundedOne{} of \resProfileIaucSIInteriorNOne{}
(\resProfileIaucSIInteriorFracOne\%, interval \resProfileIaucSIInteriorFracOneLow--\resProfileIaucSIInteriorFracOneHigh\%)
at the primary box, in \resProfileIaucSIInteriorFracHalf\% at half the box and in
\resProfileIaucSIInteriorFracTwo\% at twice the box: the clause fails at every width (H5 overall,
\resVerdictHFive; the timing clause under iAUC is met, the trace clause and the $\SI$ clause are not).
Most intervals are one-sided, so what the data constrain is a direction, and whether the other side is
bounded is decided by the box. The dependence on the box is partly a saturation of the response: the slope of
the iAUC with respect to $\SI$ at the top of the primary box is \resSatTopToBottom{} of its value at the
bottom, and at the top of the double box it is \resSatDoubleTop{} of that at the top of the primary box
(H16, exploratory). Nor is the picture an artefact of the $95\%$ level: recomputing the classification from
the stored profiles, the $\SI$ interval is bounded among interior subjects at the primary box in
\resLevelNinetySIInteriorOne\% at the $90\%$ level and \resLevelNinetyNineSIInteriorOne\% at the $99\%$ level,
and the intervals for $\ke$ and $\ka$ are bounded in \resLevelNinetyKeBoundedOne\% and
\resLevelNinetyKaBoundedOne\% of subjects at the $90\%$ level and in \resLevelNinetyNineKeBoundedOne\% and
\resLevelNinetyNineKaBoundedOne\% at the $99\%$ level (Appendix, Table~\ref{tab:levels}).

\paragraph{The replica.}
If the model were exactly true, with the subject's own noise and carbohydrate logging error of $25\%$, would
the data bound $\SI$ in the subjects where they do not? In the replica built from the real subjects
(seed zero) the $\SI$ interval is bounded in \resReplicaSIBoundedCount{} of \resReplicaSIInterior{} interior
subjects (\resReplicaSIBounded\%, interval \resReplicaSIBoundedLow--\resReplicaSIBoundedHigh\%), against
the $50\%$ that H13 asked for, and the interval for $\ke$ or $\ka$ is bounded in \resReplicaTimingMax\% of
subjects. Of the bounded $\SI$ intervals, \resReplicaCoverSI\% contain the true value, below the nominal
$95\%$ (a small count: the coverage is not estimated precisely). Verdict on H13: \resVerdictHThirteen.
Repeated over five seeds the pooled fraction is \resSeedsSIBounded\% (\resSeedsSIBoundedCount{} of
\resSeedsSIInterior{}), the fraction ranges over \resSeedsRangeLow--\resSeedsRangeHigh\% across seeds, and the
coverage of the bounded intervals is \resSeedsCoverSI\% (H17, \resVerdictHSeventeen).
\todo{Settle the wording of this paragraph from the H17 verdict: if the pooled fraction is below a quarter, say the weakness is inherent to the observable at this noise; if between a quarter and a half, say mixed.}

The replica attributes the shortfall of $\SI$ identifiability in the real data to the observable and the
design rather than to unmodelled variability, with one qualification: it is a replica of the \emph{fitted}
truths, which sit near the upper bound of the timing parameters. The synthetic study, in which the true
parameters are drawn across the box, shows where in the box recovery is possible: the $\SI$ interval is
bounded for \resSynthTercileSILow\%, \resSynthTercileSIMid\% and \resSynthTercileSIHigh\% of interior
synthetic subjects in the lower, middle and upper third of the true $\SI$ (Figure~\ref{fig:replica}b), with
coverage \resSynthCoverSI\% and a median absolute log error of the estimate of \resSynthErrorLogSI{}
($n=\resSynthN$ synthetic subjects). \todo{Interpret the tercile pattern once H19 has run.}

\begin{figure}[t]
\centering
\safefigure[width=\linewidth]{f4_replica}
\caption{Identifiability of $\SI$ when the model is true, and what the noise explains. \textbf{(a)} Subjects
with a bounded $\SI$ interval among those interior for $\SI$: real data and replica seeds. The dotted line
is the $50\%$ of H13. \textbf{(b)} Synthetic study with random truths: bounded fraction by tercile of the
true value, for $\SI$, $\ke$ and $\ka$. \textbf{(c)} Ratio of held-out iAUC error on replica data to that on
real data under the noise settings named on the axis; dashed: the $0.7$ and $0.9$ thresholds of H14.}
\label{fig:replica}
\end{figure}

\paragraph{Systematic carbohydrate error is absorbed by $\SI$.}
Scaling every logged carbohydrate by $0.75$ and by $1.33$ and refitting moves the median $\SI$ estimate by
a log ratio of \resCarbScaleLowLogRatio{} and \resCarbScaleHighLogRatio{}, against \resCarbScaleLowExpected{}
and \resCarbScaleHighExpected{} if $\SI$ absorbed the scale exactly: the same sign and a similar size. The
share of subjects with a timing rate on the upper bound is \resCarbScaleLowUpper\% and
\resCarbScaleHighUpper\%, against \resCarbScaleUnscaledUpper\% unscaled (H16, exploratory). A person whose
carbohydrate is mislogged by a constant factor will therefore be assigned a different insulin sensitivity
by the same factor, and the fit will not say so.

\paragraph{Who is recoverable.}
Whether a subject's $\SI$ interval is bounded is unrelated to the number of meals
($\rho=\resCorrMealsSIBounded$, Holm-adjusted $p=\resCorrPMealsSIBounded$), the range of carbohydrate
($\resCorrCarbRangeSIBounded$), the mean observed iAUC ($\resCorrMeanIaucSIBounded$) or the $\SI$ estimate
($\resCorrEstimateSIBounded$). Whether a timing rate is pinned on its upper bound, by contrast, tracks the
mean observed iAUC ($\rho=\resCorrMeanIaucPinned$) and a low $\SI$ estimate ($\resCorrEstimatePinned$), the
area deficit above (exploratory; Appendix, Table~\ref{tab:correlates}).

\subsection{What a richer observable buys}
\label{sec:more}

\paragraph{The window.}
The sensitivity of the iAUC to $\ke$ and $\ka$ relative to its sensitivity to $\SI$ depends on the window
(Figure~\ref{fig:window}). At ninety minutes the area is \emph{more} sensitive to timing than to $\SI$
(median ratio $\resWinNinetyLinKe$ and $\resWinNinetyLinKa$ for $\ke$ and $\ka$ in the linearized model),
at the three-hour window used here it is $\resWinOneEightyLinKe$ and $\resWinOneEightyLinKa$, and at six hours
$\resWinThreeSixtyLinKe$ and $\resWinThreeSixtyLinKa$ (H2, \resVerdictHTwo). The nonlinear engine agrees at
three hours ($\resWinOneEightyNlKe$, $\resWinOneEightyNlKa$) but not monotonically at six hours. A short
window is where timing is visible in an area; it is also where the area is least informative about $\SI$,
so the choice of window moves the problem rather than removing it.

\begin{figure}[t]
\centering
\safefigure[width=\linewidth]{f8_window}
\caption{Median (line) and interquartile range (band) over subjects of
$|\partial\,\iauc/\partial\log k|/|\partial\,\iauc/\partial\log\SI|$ for $k=\ke$ and $k=\ka$, against the
window after the meal, for the linearized and the nonlinear engine. Dashed: the thresholds of H2; dotted:
a ratio of one.}
\label{fig:window}
\end{figure}

\paragraph{The ladder of observables.}
Information about timing, measured by the condition number of the Schur-complement block after profiling out
$\SI$, rises with the richness of the observable at a common reference point
(\resCondRefIaucOne{} for iAUC, \resCondRefCentroidOne{} with the centroid, \resCondRefTraceOne{} for the
trace at the primary box; Figure~\ref{fig:fisher}b) but not at each objective's own estimate
(\resCondOwnIaucOne{}, \resCondOwnCentroidOne{}, \resCondOwnTraceOne{}), where the centroid step goes the
wrong way because those fits sit on bounds where the information is nearly singular. H3 requires the
monotone decrease at each objective's own estimate; verdict: \resVerdictHThree. The robust statement
is that the full trace carries far more timing information than the area, and the centroid adds none that is
reliable. In profile likelihood the rungs differ the same way: under the centroid no interval for $\ke$ or
$\ka$ is bounded, and under the trace \resProfileTraceKeOne\% and \resProfileTraceKaOne\% are at the primary
box and \resProfileTraceKeTwo\% and \resProfileTraceKaTwo\% at twice the box, below the $50\%$ that H5
required at the primary box (the trace clause is not met).

\begin{figure}[t]
\centering
\safefigure[width=\linewidth]{f2_fisher}
\caption{\textbf{(a)} Eigenvalues of the Fisher information in log coordinates, by observable, primary
box. \textbf{(b)} Median condition number of the timing block at a common reference point, by observable
and box.}
\label{fig:fisher}
\end{figure}

\paragraph{What is recoverable is a combination.}
Fitting in coordinates $(\log\SI,\log\tau_1,\log p)$ asks directly for the symmetric functions. The interval
for $\tau_1$ is bounded in \resCoordsBoundedTauIauc\%, \resCoordsBoundedTauCentroid\% and
\resCoordsBoundedTauTrace\% of subjects interior for it under iAUC, iAUC with the centroid and the trace;
that for $p$ in \resCoordsBoundedPIauc\%, \resCoordsBoundedPCentroid\% and \resCoordsBoundedPTrace\%
(Figure~\ref{fig:coords}). H11 required $\tau_1$ bounded in at least $70\%$ under both the centroid and
the trace, $p$ in at most $30\%$ under the centroid and at least $50\%$ under the trace; three of the four
clauses hold and the centroid clause for $\tau_1$ does not (H11, \resVerdictHEleven). In words: the
mean transit time and its companion are recoverable from the full trace, the area with its centroid gives a
minority of subjects $\tau_1$ and almost none $p$, and the area alone gives neither. A tied-rate model
$\ke=\ka$, which has only $\tau_1$ and $\SI$, predicts as well as the three-parameter model (difference in
held-out iAUC error \resTiedDiff{} mg\,dL$^{-1}$\,min, interval \resTiedDiffLow{} to \resTiedDiffHigh) but
its $\tau_1$ is bounded in only \resTiedTauBounded\% of interior subjects under iAUC with the centroid
(H12, \resVerdictHTwelve), and the coordinate fits sit on the tied boundary $p=\tau_1^2/4$ in
\resBoundaryCoords\% of cross-validation folds under iAUC and \resBoundaryCoordsTrace\% under the trace: for
many subjects the data are consistent with equal rates, not with a particular split.

\paragraph{Is the coordinate result an optimizer artefact?}
In the coordinate fits, the profile found a lower negative log-likelihood than the Adam estimate for a
large share of subjects, so the estimate was not an optimum. Replacing it with the best of six L-BFGS runs, the
gap between the two is at most \resPolishMaxGapCentroid{} (centroid) and \resPolishMaxGapTrace{} (trace) in
negative log-likelihood, below $\delta$ for \resPolishWithinCentroid\% and \resPolishWithinTrace\% of
subjects (H21, \resVerdictHTwentyOne), and on the polished estimates the bounded fractions for $\tau_1$ and
$p$ under the trace are \resPolishedBoundedTauTrace\% and \resPolishedBoundedPTrace\%. On the true-model
replica the same fits bound $\tau_1$ and $p$ under the trace in \resReplicaCoordsTauTrace\% and
\resReplicaCoordsPTrace\% of interior subjects (H18, \resVerdictHEighteen).
\todo{Interpret H18 and H21 together: whether the real-data coordinate result is reproduced by a true model with this noise.}

\begin{figure}[t]
\centering
\safefigure[width=\linewidth]{f5_coordinates}
\caption{Subjects with a bounded interval for $\tau_1$ (a) and $p$ (b) among those interior for the
parameter, by observable. Grey: the Adam estimate (real data); blue: polished by L-BFGS; green: the
true-model replica, polished.}
\label{fig:coords}
\end{figure}

\subsection{Held-out prediction cannot tell}
\label{sec:prediction}

Table~\ref{tab:prediction} and Figure~\ref{fig:prediction} give the held-out iAUC error of every cell and
the pre-specified paired differences (CGMacros, $\resCohortCgmSubjects$ subjects, five repeats of five-fold
cross-validation). Gradient fitting and exhaustive grid search of the same loss differ by
\resDiffGradThreeVsGridThree{} mg\,dL$^{-1}$\,min ($95\%$ interval \resDiffGradThreeVsGridThreeLow{} to
\resDiffGradThreeVsGridThreeHigh) with three parameters and by \resDiffGradOneVsGridOne{}
(\resDiffGradOneVsGridOneLow{} to \resDiffGradOneVsGridOneHigh) with $\SI$ only: the $90\%$ intervals lie
inside the margin of $150$ and of $90$, and H7 is met (\resVerdictHSeven). The training loss of the gradient
fit is within one percent of the grid minimum in \resGapWithinTolerance\% of folds, and a budget of five
hundred steps changes no verdict (difference \resStepsGradThreeFiveHundredVsGridThree{} with an interval of
\resStepsGradThreeFiveHundredVsGridThreeLow{} to \resStepsGradThreeFiveHundredVsGridThreeHigh). Three parameters
and $\SI$ alone differ by \resDiffGradThreeVsGradOne{} (\resDiffGradThreeVsGradOneLow{} to
\resDiffGradThreeVsGradOneHigh), equivalent at the same margins (H8, \resVerdictHEight): the two timing
parameters add nothing to held-out iAUC error. A full-trace objective does not lower it by more than the
margin: the trace objective is \resDiffGradThreeTraceVsGradThree{} (\resDiffGradThreeTraceVsGradThreeLow{} to
\resDiffGradThreeTraceVsGradThreeHigh) worse (H9, \resVerdictHNine).
What the trace objective improves is what it targets: the held-out trace error falls from
\resTraceRmseGradThree{} to \resTraceRmseGradThreeTrace{} mg\,dL$^{-1}$ and the peak-time error from
\resPeakGradThree{} to \resPeakGradThreeTrace{} min, but the personal mean, which uses no model, is at
\resTraceRmsePersonalMean{} and \resPeakPersonalMean{}, so the model adds little there either.
Against the baselines the gradient fit has lower iAUC error than the personal mean
(\resDiffGradThreeVsPersonalMean{}, \resDiffGradThreeVsPersonalMeanLow{} to \resDiffGradThreeVsPersonalMeanHigh,
better in \resWinsGradThreeVsPersonalMean{} of \resSubjectsOne{} subjects, Holm-adjusted
$p=\resPHolmGradThreeVsPersonalMean$), than a random forest (\resDiffGradThreeVsRf) and than SNPE
(\resDiffGradThreeVsSnpe); the differences from the amortized methods are differences, not equivalences, and
we do not call them equivalent.

\begin{table}[t]
\centering\small
\begin{tabular}{lcccc}
\toprule
cell & iAUC, CGMacros & iAUC, Shanghai & peak time, CGMacros & trace, CGMacros \\
 & (mg\,dL$^{-1}$\,min) & (mg\,dL$^{-1}$\,min) & (min) & (mg\,dL$^{-1}$) \\
\midrule
gradient, three parameters & 3025 & -- & 47.5 & 35.9 \\
grid search, three parameters & 3026 & -- & 48.0 & 36.3 \\
gradient, $\SI$ only & 3059 & -- & 44.4 & 33.1 \\
grid search, $\SI$ only & 3059 & -- & 44.4 & 33.1 \\
gradient, trace objective & 3057 & -- & 41.6 & 31.2 \\
random forest & 3154 & -- & 45.3 & 33.8 \\
SNPE & 3420 & -- & 43.0 & 34.8 \\
personal mean & 3257 & -- & 41.0 & 32.0 \\
population default & 3359 & -- & 46.8 & 34.8 \\
persistence & 4507 & -- & 54.6 & 46.5 \\
\bottomrule
\end{tabular}
\caption{Held-out mean absolute error (iAUC, peak time) and root mean squared error (trace) by cell; mean over subjects of the mean over five repeats of five-fold cross-validation. The Shanghai column is empty where the cell was not run.}
\label{tab:prediction}
\end{table}


\begin{figure}[t]
\centering
\safefigure[width=\linewidth]{f7_prediction}
\caption{Paired differences in held-out iAUC error between cells (first minus second; negative favours
the first), with $95\%$ bootstrap intervals over subjects. Shaded: the equivalence margin of $150$.
Blue: CGMacros; green: Shanghai.}
\label{fig:prediction}
\end{figure}

\paragraph{The same on another cohort.}
In the Shanghai cohort (\resShCvSubjects{} subjects, \resShCvRepeats{} repeats) gradient fitting and grid search
differ by \resShCvDiffGradThreeVsGridThree{} (\resShCvDiffGradThreeVsGridThreeLow{} to
\resShCvDiffGradThreeVsGridThreeHigh) and three parameters and $\SI$ alone by
\resShCvDiffGradThreeVsGradOne{} (\resShCvDiffGradThreeVsGradOneLow{} to \resShCvDiffGradThreeVsGradOneHigh),
on a personal-mean error of \resShCvMaePersonalMean{} (H20, \resVerdictHTwenty{}; post hoc).
\todo{State whether equivalence held in Shanghai and what it says about the generality of H7 and H8; if not, report the difference and say prediction is not cohort-invariant.}

\paragraph{What the remaining error is.}
A natural question is how much of the held-out error is due to noise, to optimization and to the
observable. We do not report additive shares, because no definition of the decomposition makes the terms
add. We report instead what-if noise budgets on the replica (Figure~\ref{fig:replica}c): the ratio of
the replica's held-out iAUC error to the real data's is \resReplicaRatioCgmOffCarbZero{} with no noise,
\resReplicaRatioCgmOffCarbQuarter{} with $25\%$ carbohydrate error alone, \resReplicaRatioCgmOnCarbZero{} with
CGM noise alone, \resReplicaRatioCgmOnCarbQuarter{} (\resReplicaRatioCgmOnCarbQuarterLow--\resReplicaRatioCgmOnCarbQuarterHigh)
with both, and \resReplicaRatioCgmOnCarbHalf{} with $50\%$ carbohydrate error. The primary ratio is between
the thresholds of $0.7$ and $0.9$, which the plan labelled mixed (H14). Grid search and gradient fitting give
the same error on the replica, so the optimizer is not the source. The CGM noise scale is estimated from the
residuals of the real fit, which contain model misfit as well as sensor noise, so the noise settings are an
upper bound on what sensor noise alone would explain.

\subsection{The gradient diagnostic}
\label{sec:diagnostic}

Gradient magnitude is a cheap indicator of what a parameter is worth, and we asked whether it can stand in
for the expensive analyses. With the components that push through an active bound removed, the insulin
sensitivity has the largest gradient in \resGradProjSIFirst\% of subjects (\resGradRawSIFirst\% without
the projection), and the median ratio of the timing gradient to the $\SI$ gradient is
\resGradProjRatio{} (\resGradRawRatio{} unprojected). Its ranking of the three parameters agrees with the
ranking by profile strength and by Fisher information (median Spearman \resSixMedianSpearman, with
\resSixAuc{} AUC for detecting a flat profile; H6 real-data clause). The synthetic clause, on
\resSynthN{} synthetic subjects with known truth, gives a median Spearman of \resSynthSpearman{} and an AUC of
\resSynthAuc{} (H6, \resVerdictHSix). The ordering is stable under five random initializations
($\tau=\resTenTauInit$), the two other box widths ($\resTenTauBounds$), the log parameterization
($\resTenTauParam$) and L-BFGS ($\resTenTauOptimizer$), and $\SI$ is first in \resTenSIFirst\% of the
\resTenSIFirstN{} interior subjects (H10, \resVerdictHTen). With three parameters a Spearman coefficient
takes few values, so these agreements are weaker evidence than the numbers suggest; and a diagnostic that
agrees with the profile does not make the profile unnecessary.
\todo{Rewrite this paragraph to match the H6 and H10 verdicts once the synthetic and robustness runs are in.}

\subsection{A second model class}
\label{sec:dm}

The Dalla Man meal model has a three-compartment gut with a nonlinear gastric-emptying law, so it inherits no
exact exchange symmetry from Proposition~\ref{prop:exchange}. Fitted by maximum likelihood to the same iAUC data
of the same \resDmN{} subjects, with a single-variance Gaussian noise model, its Fisher information has median
effective rank \resDmRankMedian{} of six parameters and condition number \resDmCondMedian, and the profile
interval is bounded for $k_{abs}$, $k_{max}$ and $k_{min}$ in \resDmBoundedKabs\%, \resDmBoundedKmax\% and
\resDmBoundedKmin\% of subjects, against \resDmBoundedVmx\% for its insulin-sensitivity parameter (Appendix,
Figure~\ref{fig:dm}; H22, \resVerdictHTwentyTwo). \todo{State whether the gut parameters of the second model class are also unbounded and what that says about the claim of two model classes.}

\subsection{An external-generalization check}
\label{sec:leakage}

A model that predicts well should also, one might hope, produce an $\SI$ that tracks clinical insulin
resistance. In the Shanghai cohort (\resLeakN{} subjects with HbA1c), the Spearman correlation of the
SNPE estimate of $\SI$ with HbA1c is \resLeakRawSnpe{} (\resLeakRawSnpeLow{} to \resLeakRawSnpeHigh) and that
of the gradient estimate is \resLeakRawGradient{} (\resLeakRawGradientLow{} to \resLeakRawGradientHigh);
the difference is \resLeakRawDiff{} (\resLeakRawDiffLow{} to \resLeakRawDiffHigh). With fasting glucose
partialled out the SNPE correlation falls to \resLeakPartialSnpe{} (\resLeakPartialSnpeLow{} to
\resLeakPartialSnpeHigh) and the difference between the methods to \resLeakPartialDiff{}
(\resLeakPartialDiffLow{} to \resLeakPartialDiffHigh). The SNPE estimate also correlates with its own
baseline summary feature at \resLeakSnpeFeature{} (\resLeakSnpeFeatureLow{} to \resLeakSnpeFeatureHigh).
The apparent advantage of the amortized estimate in the raw correlation is consistent with a route through
fasting glucose, not through the postprandial response; we regard this as a warning about interpreting
fitted parameters through clinical correlations, not as a result about either method (Figure~\ref{fig:leakage}).

\begin{figure}[t]
\centering
\safefigure[width=0.7\linewidth]{f9_leakage}
\caption{Spearman correlation of the estimate of $\SI$ with HbA1c in the Shanghai cohort, with $95\%$
bootstrap intervals over subjects, for the amortized (SNPE) and gradient estimates, raw and with fasting
glucose partialled out, and the correlation of the SNPE estimate with its own baseline feature.}
\label{fig:leakage}
\end{figure}

\subsection{Scorecard}
\label{sec:scorecard}

\begin{table}[t]
\centering\small
\begin{tabularx}{\linewidth}{llXl}
\toprule
& question & origin & status \\
\midrule
H1 & Exchange structure of $\ke$ and $\ka$ under symbolic analysis & plan & formulation-dependent \\
H2 & The area is insensitive to timing at long windows & plan & met \\
H3 & Timing information rises along the ladder of observables & plan & not met \\
H4 & The weakest Fisher direction is the exchange direction & plan & met \\
H5 & Bounded profile intervals: timing under iAUC and trace, $\SI$ under iAUC & plan & not met \\
H6 & The gradient diagnostic agrees with the profile likelihood & plan & partially evaluated \\
H7 & Gradient fitting and grid search predict equally well & plan & met \\
H8 & The timing parameters add nothing to held-out iAUC error & plan & met \\
H9 & A full-trace objective lowers held-out iAUC error & plan & not met \\
H10 & The diagnostic is robust to initialization, box, scale and optimizer & plan & not evaluated \\
H11 & $\tau_1$ and $p$ are identifiable under richer observables & amendment one & not met \\
H12 & A tied-rate model predicts as well and identifies $\tau_1$ & amendment one & not met \\
H13 & Replica with the model true: bounded $\SI$ interval & amendment four (post hoc) & not met \\
H14 & Replica to real held-out error ratio (noise budgets) & amendment four (post hoc) & reported (no pass/fail) \\
H15 & Replication on the Shanghai and standardized-meal cohorts & amendment four (post hoc) & met \\
H16 & Saturation, carbohydrate-scale and noise sensitivity (exploratory) & amendment four (post hoc) & reported (exploratory) \\
H17 & The replica result over five seeds & amendment five (post hoc) & not evaluated \\
H18 & $\tau_1$ and $p$ under the trace, true-model replica & amendment five (post hoc) & not evaluated \\
H19 & Synthetic recovery with random truths & amendment five (post hoc) & not evaluated \\
H20 & Prediction equivalence on the Shanghai cohort & amendment five (post hoc) & not evaluated \\
H21 & The coordinate estimates are likelihood optima & amendment five (post hoc) & not evaluated \\
H22 & A second model class (Dalla Man) & amendment five (post hoc) & not evaluated \\
\bottomrule
\end{tabularx}
\caption{Every hypothesis and its status, computed from the stored results by the rule written in the plan or the amendment (Appendix~\ref{app:register}). Primary endpoints of the plan: H1, H3, H5, H7, H9.}
\label{tab:scorecard}
\end{table}

